\section{Discussion}
\label{sec:discussion}

\paragraph{What the detectors detect.} Putting the contrasts together, the probes we tested read, in decreasing order of effect size: the prompt, the falsity of the output, and the model's knowledge. The instructed-pairs deception probe and the in-domain instructed-lie probe are dominated by the prompt; this alone produces perfect scores on the standard evaluation. Once the prompt is fixed they behave much like a truth probe, and a truth probe does the job better. A hallucination probe mixes falsity with lack of knowledge, so it rates a lie above an honest correct answer but below an honest error. None of these is a detector of belief--statement mismatch, and the feigned-ignorance contrast, where the output text is held nearly constant, shows this most directly: all are near chance, while a probe trained on that contrast reaches 0.83--0.87.

\paragraph{Relation to prior findings.} The prompt effect is the quantitative counterpart of the remark by \citet{goldowskydill2025detecting} that their probes fire on honest text related to deception, and of the finding by \citet{natarajan2026oneprobe} that prompt choice explains most of the variance in probe performance. The behaviour of the hallucination probe agrees with \citet{cheang2025really}: it tracks whether knowledge is available more than whether the output is true. The low rate of non-instructed lies and the difficulty of separating them from confusion echo \citet{hopkins2026finetuned}. What the matched design adds is the comparison none of these make: lies against falsehoods without knowledge, under one prompt.

\paragraph{Implications for evaluation.} We suggest reporting, next to the usual AUROC, (i) the AUROC between honest answers under the lie-eliciting prompt and honest answers under the control prompt, (ii) lies versus honest answers under the same prompt, and (iii) lies versus false answers on questions the model does not know. Adding the last two cells as training negatives was enough here to obtain probes that do not fire on hallucinations, although these probes still shift with the prompt and were only tested on short factual answers.

\section{Limitations}

\textbf{Scale and yield.} Both models are small (8--9B). Incentive-only prompts produced few lies, so the incentive-driven results rest on a role-play persona, a game incentive with a temperature-1 sample pool in which a quarter to a half of the ``lies'' are expected to be slips, and feigned ignorance. We did not find a setting in which these models lie about facts under a realistic incentive at a useful rate, and conclusions about strategic deception in larger models do not follow from our data. A third model (Qwen2.5-14B) was planned but not run for reasons of compute time.

\textbf{Knowledge labels.} ``Known'' and ``unknown'' are defined by sampling consistency under one neutral prompt. About 14--26\% of questions fall in between and are excluded from detector analyses. Unknown questions include consistent wrong beliefs (20\% of Llama's unknown questions and 45\% of Gemma's by a content-word heuristic), which we count as honest errors but did not analyse separately in the detector results. Wrong answers on known questions under temperature sampling include slips.

\textbf{Grading.} All labels come from a 14B open-weight judge. It agrees with GPT-4.1-mini on 92.9\% of a 9{,}473-response validation set drawn from one model; residual label noise affects all cells, and the planned API judge could not be used throughout because the API budget ran out.

\textbf{Probes and prompts.} We use one probe family (logistic regression on mean response activations at one layer), one instance of each off-the-shelf recipe, five hand-written scenarios and one question format with one-sentence answers. The instructed-pairs probe was developed for a 70B model and may behave differently there. The size of the within-prompt AUROCs varies with the layer; the layer was selected on train-split questions using the instructed scenario, which favours that scenario. The ``honest under pressure'' cell for the instructed scenario consists of cases where the model did not follow the instruction, which is itself a selected set of questions; the within-question analysis controls for the question but relies on sampling.

\textbf{Scope of ``lie''.} We only consider short factual statements that contradict a belief. Deception without false statements~\citep{thormann2026probing}, omission and exaggeration~\citep{moustafa2026beyond} are outside the design.

\section{Conclusion}

On matched questions for two open-weight models, most false statements are honest errors unless the model is told or cast to lie, and the detectors we tested respond mainly to the prompt and to falsehood in general rather than to a mismatch between belief and statement. They flag honest answers given under a deception-flavoured prompt, respond to hallucinations, and do not tell a lie from a knowledge-free falsehood or feigned from genuine ignorance. The information needed to make these distinctions is present in the activations and is recovered by probes trained against the right negatives. Evaluations of lie detectors should include those negatives.
